\minitopic{概率的多种解释}概率可以表示长期频率、客观倾向、逻辑支持或主观置信度。硬币正面概率二分之一可能指长期抛掷频率；某原子在时段内衰变的概率像客观倾向；“根据证据被告有罪的概率”表达证据支持；主体的0.7则是信心。混用解释会产生假争论。 形式认识论通常研究置信度应满足什么规范，不必先决定世界是否有客观概率。但合理置信应响应已知客观机会：若知道硬币公平，却仍给正面0.9，似乎不理性。主体概率与客观机会之间需要协调原则。 单次事件概率也不容易由频率定义。明天某候选人获胜只有一次，却可基于民调、模型和类比赋概率。频率数据是证据，模型把它转换成个案概率；转换假设必须公开。

\minitopic{概率公理与认识融贯}若主体给$A$概率0.7、给其否定概率0.6，总和超过1，信念集合内部冲突。荷兰赌论证把不融贯转成必亏赌组；准确性论证则说，非概率化置信总存在另一个概率化组合，在每种可能状态下都更接近真值\parencite{joyce1998}。 两种辩护各有前提。赌论证假设置信可用公平下注率表示，主体线性评价金钱；准确性论证需选择衡量信念离真值多远的评分规则。它们说明概率公理有强规范基础，却不表示真实人必须随时拥有无限精确数字。 区间概率允许主体用范围表达信息不足，如$P(A)\in[0.4,0.7]$。它更诚实，却使决策与更新复杂。精确概率与区间之间不是严谨和含糊之争，而是模型可操作性与不确定性表达的权衡。

\minitopic{条件化为何合理}标准贝叶斯条件化说，当主体确定学到$E$且除此之外没有新信息，新的$P'(H)=P(H\mid E)$。规则保存旧概率中关于$H$与$E$关系的判断，只重新分配到$E$成立的空间。 现实证据很少获得概率1。医生认为检验本身有10\%出错，不能把结果当确定。Jeffrey条件化允许主体只改变$E$的概率，再按加权方式更新。还有证据可能改变模型结构，使旧条件概率也需修订；机械条件化不适用。 信息顺序在理想条件化下通常不应影响最终结果，但相关证据、遗忘和模型更新会产生路径依赖。实际贝叶斯分析需记录每次证据是什么、是否独立、参数是否同时学习。

\minitopic{贝叶斯因子与证据比较}后验赔率可写为：
\[
\frac{P(H\mid E)}{P(\neg H\mid E)}=
\frac{P(H)}{P(\neg H)}\times
\frac{P(E\mid H)}{P(E\mid\neg H)}.
\]
最后一项是贝叶斯因子或似然比，表示证据把赔率放大多少。它把证据强度与先验可信度分开：强证据可显著提高一个仍不占优势的低先验假设。 假设一种症状在疾病下出现率0.8，在无病时0.2，似然比为4。若先验赔率1:99，更新后为4:99，后验仍约3.9\%。说“症状支持疾病”与“疾病现在很可能”是不同判断。 报告时同时给似然比与后验，可避免选择性叙述。只强调风险“增加四倍”可能掩盖绝对风险仍低；只强调绝对概率低又可能掩盖证据确有信息。

\minitopic{旧证据问题与新预测}一个理论常在已知证据后被提出。若$E$已经确定，简单条件化中$P(E)=1$，似乎不能再提高$H$。但相对论解释水星近日点进动等旧证据仍被视为支持。贝叶斯主义需表示“理论提出前的预期能力”、理论为何不为该数据临时拼装等元层信息。 新奇预测通常更有说服力，因为研究者较难为目标结果过拟合。不过时间上的新并非绝对要求。一个理论若从独立原则自然推出旧而精细的事实，同样有解释价值。关键是数据是否用于构造理论以及理论复杂度。 预注册、留出数据和独立复制是制度化解决方案。它们不是取代统计推理，而是让$P(E\mid H)$不因事后选择被夸大。

\minitopic{确认悖论}“所有乌鸦都是黑的”逻辑等价于“所有非黑物都不是乌鸦”。若看到一只绿苹果确认后者，按等价原则也确认前者，这就是乌鸦悖论。直觉上绿苹果对乌鸦颜色几乎无关。 贝叶斯分析可以说，绿苹果确有极微弱确认，因为它排除了一个非黑乌鸦可能，但世界中非黑非乌鸦极多，信息量近零。观察一只黑乌鸦的似然比更高。悖论提醒我们区分逻辑确认方向与实际证据强度。 另一个是基准类问题。同一个人属于多个群体，各群体频率不同。选择“同龄人”“同职业者”或“同症状者”会给不同先验。数据本身不决定基准类，因果和模型知识需要介入。

\minitopic{概率与决策}期望效用把每个结果价值乘以概率并求和。它解释低概率灾难为何可能值得预防，也说明高概率不等于应行动。设治疗成功概率0.6，成功收益大但副作用严重；是否治疗不能只看0.6超过一半。 效用不是简单金钱。健康、时间、公平和权利可能难以统一量化。决策模型可作为透明框架，显示分歧究竟在概率还是价值，而非假装产生唯一客观答案。某些道德约束也可能不允许被期望值完全抵消。 知识与决策的关系仍开放。一方说概率与效用足以，知识标签不增加决策信息；另一方认为，只有知识才可作为不再敏感分析的前提，或知识影响责任。实际高风险制度常同时使用概率阈值、证据等级与程序性知识标准。

\minitopic{群体概率与专家汇总}多个专家给出不同概率，简单平均是一种意见池，但若专家共享数据或能力差异大，平均会误导。权重可依据历史校准、领域相关性和信息独立性。权重本身也需防止事后只奖励碰巧准确者。 市场价格有时聚合分散信息，却受财富、策略、流动性和操纵影响。德尔菲法通过匿名多轮反馈减少地位压力，却可能让少数正确异议向多数收敛。没有脱离制度条件的万能聚合器。 群体还需校准：在所有给70\%的事件中约70\%发生。校准良好不等于分辨力好；永远报告基础率可能校准，却不能区分案例。Brier分数等适当评分同时奖励校准和信息性。

\minitopic{综合案例：暴雨预警}模型A给暴雨概率60\%，模型B给30\%。两者都使用同一雷达，但B另含地形数据。不能把它们当完全独立。气象员需检查历史校准、当前雷达故障、模型适用区域和共用误差。 城市是否发布红色预警还依赖错误成本：漏报可能危及生命，误报会造成停工和预警疲劳。决策阈值可以低于“最可能发生”。发布“有60\%概率”与断言“一定暴雨”承担不同认识责任。 事后没有暴雨不证明60\%预测错误。评价要看长期校准和当次输入是否正确；若模型把60\%事件预测一百次，约四十次不发生完全相容。单次结果可揭示故障，却不能单独评估概率诚实性。

\begin{tcolorbox}[breakable,colback=white,colframe=blue!30!black,title={形式分析清单},fonttitle=\bfseries]
说明概率解释；写出先验、似然与后验；检查条件概率方向和证据独立性；做先验与模型敏感性分析；区分支持度与最终概率；把概率和效用分开；用自然频率复核；报告模型限制而非只给小数。
\end{tcolorbox}
